%============================================================ 9. 결론
\section{결론}
\label{sec:conclusion}

본 연구는 자폭 무인수상정 군집에 대한 대응으로 저비용 그물을 이용한 물리적 포획을 채택하고,
이를 두 계층으로 구성한 방어 체계를 제안하였다. 두 계층은 같은 주기로 갱신되지만 판단에
걸리는 시간이 크게 달라, 느린 전략 계층의 호출만 비동기로 분리하였다. 나아가 두 계층의
기여를 720 에피소드의 시드 대응표본으로 분리 측정하였다.

기동 계층에서는 전장을 15채널 격자 지도로 래스터화하고, FiLM 조건화를 갖춘 U-Net이 낸
점수맵 위에서 픽셀을 지목하게 하여 행동공간을 규제하였다. 연속 좌표 회귀와 달리 경유점이
격자에 고정되므로 미세 진동이 구조적으로 배제되고, 유효 마스크를 통해 안전·임무 제약을
행동 분포에 직접 부과할 수 있다. 학습은 가치망 없이 후보 행동을 실제로 진행시켜 비교하는
그룹 상대 최적화로 수행하고, 한 척만 교란하여 재평가하는 반사실적 크레딧 배분으로 개별
기여를 분리하였다. 이 계층은 언어모델을 전혀 쓰지 않고 포획률을 $0.844 \to 0.886$ 으로
개선하였으며, 특히 파상 공격에서 $+0.157$ 로 컸다.

전략 계층에서는 언어모델 지휘관이 기하량이 미리 계산된 전장 상태로부터 배정을 확정하되,
조합 최적화에 의한 사후 보정을 두지 않았다. 이 설계는 성능이 아니라 측정 가능성을 위한
선택이며, 그 결과 언어모델 계층의 기여가 그대로 드러났다. 상용 클라우드 모델은 $+0.060$ 의
유의한 개선을 보였으나, 소형 로컬 모델은 $-0.192$ 로 코드 휴리스틱보다 유의하게 나빴다.

측정에서 세 가지 결과가 얻어졌다. 첫째, 두 계층의 기여는 \textbf{가산적}이다 --- 10개 지표
중 9개에서 상호작용의 신뢰구간이 0을 포함하였다. 본 논문은 이를 상승효과로 포장하지 않았다.
둘째, \textbf{포획률은 측정한 지표 중 가장 둔하다}. 효과크기 기준으로 총 선회량($-3.23$)과
포획당 그물 소모($-2.09$)의 개선이 포획률($+0.44$)보다 훨씬 크다. 다중비교 보정을 걸면 이
대조가 판정으로 바뀐다 --- 10개 지표 중 8개가 보정을 견디는데 포획률과 돌파 수만 유의성을
잃는다. 즉 제안 체계의 실질적 이득은 ``더 많이 막는다''보다 ``같은 정도로 막으면서 훨씬
적은 자원과 위험으로 해낸다''에 있다.
셋째, \textbf{언어모델 지휘관에는 능력 임계가 존재하며}, 그 성능을 설명하는 것은 배정
커버리지가 아니라 재계획 간 계획 안정성이다. 가장 낮은 성능을 보인 모델이 오히려 커버리지는
가장 높았다는 관측은, 주기적 재계획 구조에서 계획의 최적성보다 일관성이 중요하다는 설계
원리를 시사한다.

본 연구의 검증은 시뮬레이션에 한정되며, 실해역 이식과 관측 불확실성 하에서의 계획 안정성은
향후 과제로 남는다.

%============================================================ 재현성
\section*{재현성}
\addcontentsline{toc}{section}{재현성}

표 11개 가운데 7개 --- 설정 표(Table~\ref{tab:scenario})와 결과 표 6종 --- 는
\texttt{tools/make\_paper\_tables.py} 가 평가 하네스의 원자료 CSV에서 굽는다. 표 안의 숫자를
손으로 고칠 수 없다는 뜻이며, 평가를 다시 돌리면 표가 따라 바뀐다. 나머지 4개 중
Table~\ref{tab:nomenclature}, \ref{tab:metrics}, \ref{tab:conditions}는 기호·지표·조건의
\emph{정의} 표이고, 수치를 담은 수동 표는 지연 분포(Table~\ref{tab:latency_tail}) 하나뿐이며
같은 원자료에서 계산해 옮겨 적었다. 도판도 같은 원칙으로 갈린다 --- 결과 도판
(Fig.~\ref{fig:condbars}--\ref{fig:planquality})은 같은 CSV에서 생성되고, 나머지는 구조도와
개념 삽화다. 본문 문장의 수치 또한 원자료에서 계산한 것이다. 평가 실행과 표·도판 생성은
다음으로 재현된다.

% ★ \exhyphenpenalty 가 없으면 LaTeX 이 명시적 하이픈 뒤에서 줄을 끊어
%   "--backends" 가 "-- / backends" 로 갈라진다. 복사하면 안 도는 명령이 된다.
%   줄은 좁은 단에 맞춰 손으로 끊고, 이어지는 줄은 들여쓴다.
\begin{quote}\ttfamily\footnotesize\raggedright
\exhyphenpenalty=10000 \hyphenpenalty=10000
\# 조건행렬 전체 수집 (30 시드)\\
python run\_eval.py --seeds 30 \textbackslash\\
\hspace*{1.4em}--backends ollama,gemini\\[3pt]
\# 표 생성 (원자료 CSV $\to$ tables/*.tex)\\
python tools/make\_paper\_tables.py\\[3pt]
\# 논문 결과 도판 (Fig.~\ref{fig:condbars}--\ref{fig:planquality})\\
python -m boatattack\_sim.eval.paper\_figs \textbackslash\\
\hspace*{1.4em}--csv results/eval\_merged\\[3pt]
\# 전면부·기호 규약 검사\\
python tools/check\_paper\_meta.py
\end{quote}

\noindent
정책 체크포인트와 환경 설정, 그리고 계획 호출의 계측값(지연·폴백·커버리지·churn·교차·근거 길이)을 보존하였다. 근거 원문은 저장하지 않았다.
